% Part II, introduction and related work (included by main_en.tex).

\section{Introduction}\label{sec:intro}

Reinforcement learning with verifiable rewards (RLVR) trains language models on problems whose answers can be checked.
GRPO~\cite{shao2024deepseekmath,deepseek2025r1} and its variants~\cite{liu2025understanding,liu2026tricks} are the most
common form. The verifier is exact but coarse. It says whether the final answer is right, not whether the reasoning was
sound, and it gives no signal in groups where every sample is right or every sample is wrong. A judge model can grade
the reasoning itself. Process supervision~\cite{lightman2024verify,wang2024mathshepherd} and generative
verifiers~\cite{zhang2024genverifier} are two forms of this idea. A judge brings a cost, in memory and time, and a risk:
a learned reward can be exploited~\cite{gao2023scaling,skalse2022defining}.

\kev{}~\cite{kev4b} is a judge that is unusually cheap to add to RLVR on a Qwen3.5-4B policy. It is a LoRA adapter and
a pointer head on the same base model, Qwen3.5-4B-Base, so the policy and the judge can share one copy of the backbone.
Policy LoRA and judge LoRA sit side by side, and the judge reads its answer in one forward pass instead of generating
text. This paper measures what that design saves and what the judge adds to GRPO. We ask three questions:
\begin{enumerate}
\item How much GPU memory does the shared backbone save over loading the judge as a second model, measured in the
  same process with the same training step (\S\ref{sec:res-shared})?
\item How well does \kev{} tell correct from incorrect solutions of the policy, compared with an AR judge of the
  same size and family, and with that judge reasoning before it answers? How easily is each judge moved by edits that
  should not move it (\S\ref{sec:res-study})?
\item Does adding \kev{}'s judgment to the verifier's reward improve GRPO on MATH, and can \kev{} replace the
  verifier (\S\ref{sec:res-grpo})?
\end{enumerate}

\paragraph{Correction.} Versions 1--3 of this paper reported a leaderboard, learning curves and memory figures that were
not produced by experiments. The leaderboard values were constants in a script. The memory of ``two separate models''
was set to 2.04 times the measured single-backbone memory, and the memory of ``two adapters'' to that measurement plus
0.12\,GB (see the notice above and \texttt{CORRECTION.md}). This version measures all of them. The training runs use
fewer problems, steps and benchmarks than v1--v3 claimed, which is what we could run and check. Every number is
computed from the logs in the repository.

\section{Related work}\label{sec:related}

\paragraph{RLVR and GRPO.} GRPO~\cite{shao2024deepseekmath} replaces PPO's value function with a group baseline over
several samples of the same prompt. DeepSeek-R1~\cite{deepseek2025r1} made it the standard recipe for reasoning. Later
analyses showed that details matter: length normalisation and the per-group standard
deviation~\cite{liu2025understanding}, and the choice of advantage scale~\cite{liu2026tricks}. We use the group mean with
the batch-level scale that \citet{liu2026tricks} recommend, a token-level loss, and no KL term. Modern RL stacks sample
with an inference engine and train with another framework. The two compute slightly different probabilities for the
same weights, and truncated importance sampling corrects for that~\cite{yao2025offpolicy}.

\paragraph{Judges and reward models for reasoning.} Outcome and process reward models score solutions or
steps~\cite{lightman2024verify,wang2024mathshepherd}. Generative verifiers~\cite{zhang2024genverifier} read the
probability of ``yes'' from a language model asked whether a solution is correct, optionally after it reasons. LLM
judges are known to be sensitive to length, style and assertions of correctness~\cite{zheng2023judging}. Reward-model
over-optimisation grows with how hard the policy is pushed~\cite{gao2023scaling}, and RL policies have learned to
exploit their reward signals in many settings~\cite{skalse2022defining,denison2024sycophancy,baker2025monitoring}.

\paragraph{Sharing weights between policy and reward.} Multi-LoRA serving~\cite{hu2022lora,kwon2023vllm} lets several
adapters share one base model at inference time. Here the same mechanism is used inside the trainer, for a policy
adapter that is being trained and a judge adapter that is frozen. The condition that makes this exact is that training
never changes the base weights the judge was trained on.
